\emph{(1)} In $\Cl(1,3)$ with $e_0^2 = 1$, $e_i^2 = -1$: $B_i^2 = e_0e_ie_0e_i = -e_0^2e_i^2 = 1$, and for $i \neq j$, $B_iB_j = e_0e_ie_0e_j = -e_ie_j$, so $B_iB_j + B_jB_i = -(e_ie_j + e_je_i) = 0$. Hence $B_w^2 = \sum_i p_i^2 B_i^2 + \sum_{i<j}p_ip_j(B_iB_j + B_jB_i) = \sum_i p_i^2$, a scalar; checked exactly at two circle points and two disk points. With $\sum_i p_i = 1$ (the first Kasner relation, \claimref{dark-energy.frozen-point-datum-kasner-wall}(ii)), $B_w^2 = 1$ is the second relation.

\emph{(2)} $1 = (\sum_i p_i)^2 = \sum_i p_i^2 + 2\sum_{i<j}p_ip_j$, so $\sum_i p_i^2 = 1 - 2\sum_{i<j}p_ip_j$. If every $p_i > 0$ the cross sum is positive and $\sum_i p_i^2 < 1$. On the circle the cross sum vanishes, so the exponents cannot all be positive: some $p_i \leq 0$. If $p_1 = 0$ on the circle then $p_2 + p_3 = 1 = p_2^2 + p_3^2$, so $p_2p_3 = 0$ and $p$ is a permutation of $(1,0,0)$. Checked: the rational parametrisation $p(u) = (-u, 1+u, u(1+u))/(1+u+u^2)$ satisfies both relations and has $p_1 < 0$ for $u > 0$; the zero-entry case on a grid; and every positive triple $(a,b,c)/12$ (55 of them) has $\sum p_i^2 < 1$, maximum $17/24$, with the cross-sum identity.

\emph{(3)} \claimref{dark-energy.frozen-point-datum-kasner-wall}(ii) gives, in its convention (de Sitter has $\mathrm{Ric} = \Lambda g$), $\tau^2R_{00} = \sum p_i^2 - \sum p_i$ and $\tau^2 R_{ii} = -p_i(\sum p_j - 1)$, the rest vanishing, on the Kasner ansatz with constant exponents. In this convention a vacuum energy $\rho g_{ab}$ with $\rho = \Lambda/8\pi G$ has trace-reversed stress $-\rho g_{ab}$ and must produce $\mathrm{Ric} = \Lambda g$, so the Einstein equation reads $R_{ab} = -8\pi G\,(T_{ab} - \tfrac12 T g_{ab})$. For a massless scalar, $T_{ab} - \tfrac12Tg_{ab} = \partial_a\phi\,\partial_b\phi$ in any signature, and with $\phi = q\log\tau$ this is $q^2\tau^{-2}$ in the $00$ slot and zero elsewhere. So $R_{ii} = 0$, which for $p \neq 0$ is $\sum p_i = 1$, and $\sum p_i^2 - 1 = -8\pi G q^2$. The wave equation $\Box\phi = \tau^{-\sum p_i}\partial_\tau(\tau^{\sum p_i}\partial_\tau\phi) = \tau^{-1}\partial_\tau(q) = 0$ holds exactly when $\sum p_i = 1$. The off-diagonal $0i$ equations are those of \claimref{synthesis.diagonal-kasner-momentum-constraint} with the scalar's contribution $\partial_0\phi\,\partial_i\phi = 0$, so they are unchanged (constant exponents). Then $B_w^2 = \sum p_i^2 = 1 - 8\pi G q^2$ by (1). The isotropic stiff point $(\tfrac13,\tfrac13,\tfrac13)$, the stiff-fluid Friedmann expansion $a \propto \tau^{1/3}$, has $8\pi Gq^2 = \tfrac23$. Checked: the relations at three points, including that the residue square equals $1 - 8\pi G q^2$ computed from $R_{00}$.

\emph{The last sentence.} Step (2) of \claimref{dark-energy.forcing-half-schwarzschild-wall-residue-sense} states that the residues satisfy the wall's field equations and that the second Kasner relation is the unit-boost property. By (3) the field equations at a wall with a scalar of nonzero log-rate give a residue of norm squared below one, so the equivalence holds only for $q = 0$, the vacuum case, and then (2) gives an exponent $p_i \leq 0$.

\emph{What is argued, not proved.} The dynamical reading in the technical level uses literature recalled from memory and not checked here: in vacuum, a Kasner regime with an exponent $p_1 < 0$ is destabilised by the spatial-curvature term growing like $\tau^{4p_1}$ unless its structure constant vanishes (Belinski, Khalatnikov and Lifshitz, 1970; for Bianchi IX, Ringström, 2001, proved convergence to the oscillatory attractor), and the Fournodavlos-Luk class (2020) is built so that this term vanishes; with a massless scalar, solutions with all exponents positive are quiescent and exist in a full-dimensional family (Andersson and Rendall, 2001). These place the obstruction: a generic extension of the forcing half beyond the diagonal vacuum class either has no limiting residue or has a residue that is not a unit boost.
